\section{FFTDF: reciprocal-space Coulomb and exchange derivatives}
\label{sec:fftdf}

\subsection{Uniform-grid and plane-wave representations}

FFTDF evaluates AO-pair densities on a uniform real-space grid and transforms
them to reciprocal space with an FFT.  Schematically, the total-density
Fourier coefficient generated by the density matrices is
\begin{equation}
  \rho_D(\vG)
  =\frac{1}{N_k}\sum_{\vk\mu\nu}
   D_{\mu\nu}^{\vk}
   \int_\Omega e^{-i\vG\cdot\mathbf r}
   \phi_{\nu\vk}^{*}(\mathbf r)
   \phi_{\mu\vk}(\mathbf r)\dd\mathbf r.
  \label{eq:fftdf-density-fourier}
\end{equation}
For a neutral periodic system, excluding the $\vG=0$ component, the Hartree
energy is
\begin{equation}
  E_{\mathrm H}
  =\frac{1}{2\Omega}\sum_{\vG\ne0}
   \frac{4\pi}{G^2}|\rho_D(\vG)|^2.
  \label{eq:fftdf-hartree-energy}
\end{equation}
Inverse-transforming the Hartree potential and integrating it with the AO
pairs produces $J^{\vk}$.

\subsection{Nuclear derivatives at fixed mesh}

At fixed cell and fixed integer mesh, neither the grid points nor the $\vG$ set
change under a nuclear displacement.  The nuclear derivative of the Coulomb
kernel $4\pi/G^2$ therefore vanishes, leaving only the derivative of the AO-pair
density,
\begin{equation}
  E_{\mathrm H}^{Ax}
  =\frac{1}{\Omega}\operatorname{Re}
   \sum_{\vG\ne0}\frac{4\pi}{G^2}
   \rho_D(\vG)^*\rho_D^{Ax}(\vG).
  \label{eq:fftdf-hartree-derivative}
\end{equation}
$\rho_D^{Ax}$ contains the AO row and column derivatives for atom $A$ and the
same periodic-image and Bloch-phase conventions as the central density.
Constructing an AO-row matrix followed by Hermitian reconstruction yields the
full $J^{Ax}$ entering Eq.~\eqref{eq:hartree-gradient}.

\subsection{Exact exchange and momentum transfer}

Source HF exchange uses the Coulomb kernel $4\pi/|\vG+\vq|^2$ for each pair
momentum transfer
\begin{equation}
  \vq=\vk'-\vk.
  \label{eq:fftdf-momentum-transfer}
\end{equation}
Symbolically,
\begin{equation}
  K_{\mu\nu}^{\vk}
  =\frac{1}{N_k}\sum_{\vk'}\sum_{\lambda\sigma}
   (\mu\vk,\lambda\vk'|\sigma\vk',\nu\vk)
   D_{\lambda\sigma}^{\vk'}.
  \label{eq:fftdf-exchange}
\end{equation}
The nuclear derivative applies the product rule to whichever of the four AO
centers moves.  The same FFT exchange provider must be applied to transition
densities in the source HF Hessian action.

\subsection{\texorpdfstring{$G=0$, $q=0$}{G=0, q=0}, and the Ewald exchange divergence}

The Hartree $\vG=0$ mode is omitted or corrected consistently with the
neutralizing convention and GTH ionic charges.  Exchange is singular at
$\vq=0$, $\vG=0$, and PySCF's \texttt{exxdiv='ewald'} convention adds an
overlap-dependent correction.  Because the overlap changes under nuclear
motion, this correction also requires a derivative of the form
\begin{equation}
  K_{\mathrm{Ewald}}^{Ax}
  \propto S^{Ax}DS+SDS^{Ax}.
  \label{eq:fftdf-ewald-exchange-derivative}
\end{equation}
Confusing ``Gamma'' with ``$\vq=0$'' loses the singular correction for the
diagonal pairs of a sampled-k calculation.

\subsection{Cutoff and mesh}

The plane-wave kinetic cutoff is approximately
\begin{equation}
  E_{\mathrm{cut}}=\frac{G_{\max}^2}{2}.
  \label{eq:cutoff-definition}
\end{equation}
Considering the Nyquist wavevector along each direction of an orthorhombic
cell, the mesh size scales with the cell length and with
$\sqrt{E_{\mathrm{cut}}}$.  PySCF couples \texttt{cell.ke\_cutoff} and
\texttt{cell.mesh}, but integer rounding means that a small cutoff change can
change the mesh discontinuously.

Both $\vR+h$ and $\vR-h$ in a full-energy finite difference must use the same
lattice and the same integer mesh.  Re-estimating the mesh from the cutoff at
each displaced geometry mixes a change of numerical representation into the
finite difference.  Metadata must therefore record both the cutoff and the
resolved mesh.

\subsection{Terms owned by FFTDF in the gradient}

A consistent FFTDF provider supplies
\begin{enumerate}
  \item the central source-HF $J[D]$ and $K[D]$,
  \item $\delta J$ and $\delta K$ for transition densities in AX,
  \item the full fixed-density $J^{Ax}[D]$ and $K^{Ax}[D]$,
  \item the $\vq=0$ Ewald exchange-divergence derivative, and
  \item the explicit target-Hartree derivative.
\end{enumerate}
Target XC is owned by the NumInt/MGGA grid provider, while the GTH
one-electron derivative is owned by a separate one-electron provider.  Keeping
these responsibilities separate avoids component double counting even when
the same mesh is used.

\subsection{CPU and GPU ownership}

The CPU route uses the PySCF periodic FFTDF central value and project-local
derivative matrices.  The GPU route uses GPU4PySCF FFTs, AO grids, and CuPy
contractions.  Cell and lattice objects, small pseudopotential tables, and
final reporting may remain on the host, but AO/grid FFTs, J/K, AX, the
Z-vector, and final electronic assembly must be device operations.  The
capability report discloses host boundaries and forbids silent CPU fallback.

FFTDF validation does not end when a finite difference agrees at one
$E_{\mathrm{cut}}$.  Energy and gradient must converge along a resolved mesh
sequence, Gamma must agree with an equivalent one-k-point representation, and
Eq.~\eqref{eq:translation-sum-rule} must be satisfied.
